% GENERATED FILE. Edit canonical lessons, then run npm run generate:books.
% canonical-source-hash: fnv1a64:3ac7767ac277f47e
% canonical-lessons: GU-C241-khilvu, GU-C241-karmavu, GU-C241-mul, GU-C241-ukhedvu, GU-C241-marvu

\chapter{Flowers that bloomed and died}
\label{ch:gu-flowers-that-bloomed-and-died}
\hlchaptermodality{241}

\begin{chapteropening}
\textbf{By the end of this chapter:} \emph{I can say how my plants grew, bloomed and died.}

Complete the last lesson of chapter 241: I can say how my plants grew, bloomed and died.
\end{chapteropening}

\section[\texorpdfstring{\gu{ખીલવું} (khīlvũ)}{ખીલવું (khīlvũ)}]{\gu{ખીલવું} (khīlvũ) — to bloom, to blossom}

\label{lesson:GU-C241-khilvu}



\begin{quote}
Before the new one: say the Gujarati for to irrigate, then the Gujarati for rose.
\end{quote}

\subsection*{You'll want to know: \gu{ખીલવું}}
\textbf{\gu{ખીલવું}} — \emph{khīlvũ} — \textquotedblleft{}to bloom, to blossom\textquotedblright{}.

\begin{quote}
\textbf{\gu{ગઈકાલે} \gu{ગુલાબ} \gu{ખીલ્યું}.} — \emph{gaīkāle gulāb khīlyũ} — \textquotedblleft{}Yesterday the rose bloomed.\textquotedblright{}
\end{quote}

\begin{cousinweb}
\textbf{\gu{ખીલવું}} (\emph{khīlvũ}) comes from \emph{*khill-}, \textquotedblleft{}be open, expand\textquotedblright{}, one of the reconstructed forms that CDIAL gathers under \emph{*khiḍ}; it glosses the Gujarati verb \textquotedblleft{}to bloom, expand\textquotedblright{} (CDIAL 3882).
\end{cousinweb}

\begin{grammarlens}[title={What you've built}]
One more word for everyday life at A2.
\end{grammarlens}

\subsection*{Your turn}
\begin{itemize}
\raggedright
  \item \emph{Say it:} \emph{khīlvũ}
  \item \emph{Say it:} \emph{khīlvũ}, once more
  \item \emph{Say it:} \emph{gaīkāle gulāb khīlyũ} — Yesterday the rose bloomed.
  \item \emph{Recall:} say the Gujarati for to irrigate, then the Gujarati for rose, then say \emph{khīlvũ} again
\end{itemize}

\begin{tcolorbox}[breakable,colback=teal!4,colframe=teal!35!black,title={Before you move on}]
What does \gu{ખીલવું} mean? (\textquotedblleft{}To bloom, to blossom\textquotedblright{}.) Say it once more.
\end{tcolorbox}
\section[\texorpdfstring{\gu{કરમાવું} (karmāvũ)}{કરમાવું (karmāvũ)}]{\gu{કરમાવું} (karmāvũ) — to wilt, to wither}

\label{lesson:GU-C241-karmavu}



\begin{quote}
Before the new one: say the Gujarati for rose, then the Gujarati for to bloom.
\end{quote}

\subsection*{You'll want to know: \gu{કરમાવું}}
\textbf{\gu{કરમાવું}} — \emph{karmāvũ} — \textquotedblleft{}to wilt, to wither\textquotedblright{}.

\begin{quote}
\textbf{\gu{તડકામાં} \gu{ફૂલ} \gu{કરમાયું}.} — \emph{taḍkāmā̃ phūl karmāyũ} — \textquotedblleft{}The flower wilted in the sun.\textquotedblright{}
\end{quote}

\begin{cousinweb}
\textbf{\gu{કરમાવું}} (\emph{karmāvũ}) comes from \emph{*kurumm-}, \textquotedblleft{}wither\textquotedblright{}, a reconstructed form that CDIAL lists beside \emph{*komh}; it glosses the Gujarati verb \textquotedblleft{}wither\textquotedblright{} (CDIAL 3524).
\end{cousinweb}

\begin{grammarlens}[title={What you've built}]
One more word for everyday life at A2.
\end{grammarlens}

\subsection*{Your turn}
\begin{itemize}
\raggedright
  \item \emph{Say it:} \emph{karmāvũ}
  \item \emph{Say it:} \emph{karmāvũ}, once more
  \item \emph{Say it:} \emph{taḍkāmā̃ phūl karmāyũ} — The flower wilted in the sun.
  \item \emph{Recall:} say the Gujarati for rose, then the Gujarati for to bloom, then say \emph{karmāvũ} again
\end{itemize}

\begin{tcolorbox}[breakable,colback=teal!4,colframe=teal!35!black,title={Before you move on}]
What does \gu{કરમાવું} mean? (\textquotedblleft{}To wilt, to wither\textquotedblright{}.) Say it once more.
\end{tcolorbox}
\section[\texorpdfstring{\gu{મૂળ} (mūḷ)}{મૂળ (mūḷ)}]{\gu{મૂળ} (mūḷ) — root}

\label{lesson:GU-C241-mul}



\begin{quote}
Before the new one: say the Gujarati for to bloom, then the Gujarati for to wilt.
\end{quote}

\subsection*{You'll want to know: \gu{મૂળ}}
\textbf{\gu{મૂળ}} — \emph{mūḷ} — \textquotedblleft{}root\textquotedblright{}.

\begin{quote}
\textbf{\gu{છોડનું} \gu{મૂળ} \gu{સૂકું} \gu{હતું}.} — \emph{choḍnũ mūḷ sūkũ hatũ} — \textquotedblleft{}The plant's root was dry.\textquotedblright{}
\end{quote}

\begin{cousinweb}
\textbf{\gu{મૂળ}} (\emph{mūḷ}) is a tadbhava. It comes from Sanskrit \emph{mūla}, \textquotedblleft{}root\textquotedblright{}, and CDIAL glosses the Gujarati word \textquotedblleft{}root of tree, origin\textquotedblright{} (CDIAL 10250).
\end{cousinweb}

\begin{grammarlens}[title={What you've built}]
One more word for everyday life at A2.
\end{grammarlens}

\subsection*{Your turn}
\begin{itemize}
\raggedright
  \item \emph{Say it:} \emph{mūḷ}
  \item \emph{Say it:} \emph{mūḷ}, once more
  \item \emph{Say it:} \emph{choḍnũ mūḷ sūkũ hatũ} — The plant's root was dry.
  \item \emph{Recall:} say the Gujarati for to bloom, then the Gujarati for to wilt, then say \emph{mūḷ} again
\end{itemize}

\begin{tcolorbox}[breakable,colback=teal!4,colframe=teal!35!black,title={Before you move on}]
What does \gu{મૂળ} mean? (\textquotedblleft{}Root\textquotedblright{}.) Say it once more.
\end{tcolorbox}
\section[\texorpdfstring{\gu{ઉખેડવું} (ukheḍvũ)}{ઉખેડવું (ukheḍvũ)}]{\gu{ઉખેડવું} (ukheḍvũ) — to uproot, to pull up}

\label{lesson:GU-C241-ukhedvu}



\begin{quote}
Before the new one: say the Gujarati for to wilt, then the Gujarati for root.
\end{quote}

\subsection*{You'll want to know: \gu{ઉખેડવું}}
\textbf{\gu{ઉખેડવું}} — \emph{ukheḍvũ} — \textquotedblleft{}to uproot, to pull up\textquotedblright{}.

\begin{quote}
\textbf{\gu{મેં} \gu{સૂકો} \gu{છોડ} \gu{ઉખેડ્યો}.} — \emph{mẽ sūko choḍ ukheḍyo} — \textquotedblleft{}I pulled up the dry plant.\textquotedblright{}
\end{quote}

\begin{cousinweb}
\textbf{\gu{ઉખેડવું}} (\emph{ukheḍvũ}) is a tadbhava. It comes from \emph{*utskarati}, \textquotedblleft{}digs out\textquotedblright{}, a reconstructed form; CDIAL glosses the Gujarati verb \textquotedblleft{}to root up\textquotedblright{} (CDIAL 1895).
\end{cousinweb}

\begin{grammarlens}[title={What you've built}]
One more word for everyday life at A2.
\end{grammarlens}

\subsection*{Your turn}
\begin{itemize}
\raggedright
  \item \emph{Say it:} \emph{ukheḍvũ}
  \item \emph{Say it:} \emph{ukheḍvũ}, once more
  \item \emph{Say it:} \emph{mẽ sūko choḍ ukheḍyo} — I pulled up the dry plant.
  \item \emph{Recall:} say the Gujarati for to wilt, then the Gujarati for root, then say \emph{ukheḍvũ} again
\end{itemize}

\begin{tcolorbox}[breakable,colback=teal!4,colframe=teal!35!black,title={Before you move on}]
What does \gu{ઉખેડવું} mean? (\textquotedblleft{}To uproot, to pull up\textquotedblright{}.) Say it once more.
\end{tcolorbox}
\section[\texorpdfstring{\gu{મરવું} (marvũ)}{મરવું (marvũ)}]{\gu{મરવું} (marvũ) — to die}

\label{lesson:GU-C241-marvu}



\begin{quote}
Before the new one: say the Gujarati for root, then the Gujarati for to uproot.
\end{quote}

\subsection*{You'll want to know: \gu{મરવું}}
\textbf{\gu{મરવું}} — \emph{marvũ} — \textquotedblleft{}to die\textquotedblright{}.

\begin{quote}
\textbf{\gu{ગરમીમાં} \gu{છોડ} \gu{મર્યો}.} — \emph{garmīmā̃ choḍ maryo} — \textquotedblleft{}The plant died in the heat.\textquotedblright{}
\end{quote}

\begin{cousinweb}
\textbf{\gu{મરવું}} (\emph{marvũ}) is a tadbhava, a word inherited from Sanskrit and reshaped by sound change. It comes from \emph{marate}, \textquotedblleft{}will die\textquotedblright{} (CDIAL 9871).
\end{cousinweb}

\begin{grammarlens}[title={What you've built}]
One more word for everyday life at A2.
\end{grammarlens}

\subsection*{Your turn}
\begin{itemize}
\raggedright
  \item \emph{Say it:} \emph{marvũ}
  \item \emph{Say it:} \emph{marvũ}, once more
  \item \emph{Say it:} \emph{garmīmā̃ choḍ maryo} — The plant died in the heat.
  \item \emph{Recall:} say the Gujarati for root, then the Gujarati for to uproot, then say \emph{marvũ} again
\end{itemize}

\begin{tcolorbox}[breakable,colback=teal!4,colframe=teal!35!black,title={Before you move on}]
What does \gu{મરવું} mean? (\textquotedblleft{}To die\textquotedblright{}.) Say it once more.
\end{tcolorbox}
